% -*- Mode: LaTeX; Package: CLIM-USER -*-

\chapter {Graph Formatting}
\label {graph-formatting}

CLIM provides a mechanism for arranging arbitrary output in a graph.  The
following code produces the graph shown in Figure~\ref{graph-example}.

\begin{verbatim}
(defun graph-test (stream &optional (orientation :horizontal)) 
  (fresh-line stream)
  (macrolet ((make-node (&key name children)
               `(list* ,name ,children)))
    (flet ((node-name (node)
             (car node))
           (node-children (node)
             (cdr node)))
      (let* ((2a (make-node :name "2A"))
             (2b (make-node :name "2B"))
             (2c (make-node :name "2C"))
             (1a (make-node :name "1A" :children (list 2a 2b)))
             (1b (make-node :name "1B" :children (list 2b 2c)))
             (root (make-node :name "0" :children (list 1a 1b))))
        (format-graph-from-roots
          (list root)
          #'(lambda (node s)
              (write-string (node-name node) s))
          #'node-children
          :orientation orientation
          :stream stream)))))
\end{verbatim}

\begin{figure}
\centerline{\includegraphics{graph-example}}
\caption{\label{graph-example} Example of graph formatting.}
\end{figure}


\section {Graph Formatting Functions}

\Defun {format-graph-from-roots} {root-objects object-printer inferior-producer
                                  \key stream
                                       orientation cutoff-depth
                                       merge-duplicates duplicate-key duplicate-test
                                       generation-separation within-generation-separation
                                       center-nodes
                                       arc-drawer arc-drawing-options
                                       graph-type (move-cursor \cl{t})}

Draws a graph whose roots are specified by the sequence \arg{root-objects}.  The
nodes of the graph are displayed by calling the function \arg{object-printer},
which takes two arguments, the node to display and a stream.
\arg{inferior-producer} is a function of one argument that is called on each
node to produce a sequence of inferiors (or \cl{nil} if there are none).  Both
\arg{object-printer} and \arg{inferior-producer} have dynamic extent.

The output from graph formatting takes place in a normalized +$y$-downward
coordinate system.  The graph is placed so that the upper left corner of its
bounding rectangle is at the current text cursor position of \arg{stream}.  If
the boolean \arg{move-cursor} is \term{true} (the default), then the text cursor
will be moved so that it immediately follows the lower right corner of the
graph.

The returned value is the output record corresponding to the graph.

\arg{stream} is an output recording stream to which output will be done.  It
defaults to \cl{*standard-output*}.

\arg{orientation} may be either \cl{:horizontal} (the default) or
\cl{:vertical}.  It specifies which way the graph is oriented.  CLIM
implementations are permitted to extend the values of \arg{orientation}, for
example, adding \cl{:right} or \cl{:left} to distinguish between left-to-right
or right-to-left layouts.

\arg{cutoff-depth} specifies the maximum depth of the graph.  It defaults to
\cl{nil}, meaning that there is no cutoff depth.  Otherwise it must be an
integer, meaning that no nodes deeper than \arg{cutoff-depth} will be formatted
or displayed.

If the boolean \arg{merge-duplicates} is \term{true}, then duplicate objects in
the graph will share the same node in the display of the graph.  That is, when
\arg{merge-duplicates} is \term{true}, the resulting graph will be a tree.  If
\arg{merge-duplicates} is \term{false} (the default), then duplicate objects
will be displayed in separate nodes.  \arg{duplicate-key} is a function of one
argument that is used to extract the node object component used for duplicate
comparison; the default is \cl{identity}.  \arg{duplicate-test} is a function of
two arguments that is used to compare two objects to see if they are duplicates;
the default is \cl{eql}.  \arg{duplicate-key} and \arg{duplicate-test} have
dynamic extent.

\arg{generation-separation} is the amount of space to leave between successive
generations of the graph; the default should be chosen so that the resulting
graph is visually pleasing.  \arg{within-generation-separation} is the amount of
space to leave between nodes in the same generation of the graph; the default
should be chosen so that the resulting graph is visually pleasing.
\arg{generation-separation} and \arg{within-generation-separation} are specified
in the same way as the \arg{inter-row-spacing} argument to
\cl{formatting-table}.

When \arg{center-nodes} is \term{true}, each node of the graph is centered with
respect to the widest node in the same generation.  The default is \term{false}.

\arg{arc-drawer} is a function of seven positional and some unspecified keyword
arguments that is responsible for drawing the arcs from one node to another; it
has dynamic extent.  The positional arguments are the stream, the ``from''
node's object, the ``to'' node's object, the ``from'' $x$ and $y$ position, and
the ``to'' $x$ and $y$ position.  The keyword arguments gotten from
\arg{arc-drawing-options} are typically line drawing options, such as for
\cl{draw-line*}.  If \arg{arc-drawer} is unsupplied, the default behavior is to
draw a thin line from the ``from'' node to the ``to'' node using
\cl{draw-line*}.

\arg{graph-type} is a keyword that specifies the type of graph to draw.  All
CLIM implementations must support graphs of type \cl{:tree},
\cl{:directed-graph} (and its synonym \cl{:digraph}), and
\cl{:directed-acyclic-graph} (and its synonym \cl{:dag}).  \arg{graph-type}
defaults to \cl{:digraph} when \arg{merge-duplicates} is \term{true}, otherwise
it defaults to \cl{:tree}.  Typically, different graph types will use different
output record classes and layout engines to lay out the graph.  However, it is
permissible for all of the required graph types to use exactly the same graph
layout engine.


\section {The Graph Formatting Protocols}

Graph formatting is implemented on top of the basic output recording protocol,
using \cl{with-new-output-record} to specify the appropriate type of output
record.  For example, \cl{format-graph-from-roots} first collects all the output
that belongs in the graph into a collection of graph node output records by
calling \cl{generate-graph-nodes}.  All of the graph node output records are
descendents of a single graph output record.  During this phase,
\cl{stream-drawing-p} is bound to \cl{nil} and \cl{stream-recording-p} is bound
to \cl{t}.  When all the output has been generated, the graph layout code
(\cl{layout-graph-nodes} and \cl{layout-graph-edges}) is called to compute the
graph layout.  Finally, the graph output record is positioned on the stream at
the current text cursor position and then displayed by calling \cl{replay} on
the graph output record.


\Defprotoclass {graph-output-record}

The protocol class that represents a graph; a subclass of \cl{output-record}.
\IfYouWantClass {a} {graph output record} {graph-output-record}

\Defpredicate {graph-output-record-p} {object}

Returns \term{true} if \arg{object} is a \term{graph output record}, otherwise
returns \term{false}.

\Defclass {standard-graph-output-record}

The instantiable class of output record that represents a graph.  Its children
will be a sequence graph nodes.  This is a subclass of \cl{graph-output-record}.

\definitarg {:orientation}
\definitarg {:center-nodes}
\definitarg {:cutoff-depth}
\definitarg {:merge-duplicates}
\definitarg {:generation-separation}
\definitarg {:within-generation-separation}
\Definitarg {:hash-table}

All the graph output record must handle these seven initargs, which are used to
specify, respectively, the orientation, node centering, cutoff depth, merge
duplicates, generation and within-generation spacing, and the node hash table of
a graph output record.


\Defmacro {define-graph-type} {graph-type class}

Defines a new graph type named by the symbol \arg{graph-type} that is
implemented by the class \arg{class}.  \arg{class} must be a subclass of
\cl{graph-output-record}.  Neither of the arguments is evaluated.

All CLIM implementations must support graphs of type \cl{:tree},
\cl{:directed-graph} (and its synonym \cl{:digraph}), and
\cl{:directed-acyclic-graph} (and its synonym \cl{:dag}).


\Defgeneric {graph-root-nodes} {graph-record}

Returns a sequence of the graph node output records corresponding to the root
objects for the graph output record \arg{graph-record}.

\Defgeneric {(setf graph-root-nodes)} {roots graph-record}

Sets the root nodes of \arg{graph-record} to \arg{roots}.


\Defgeneric {generate-graph-nodes} {graph-record stream
                                    root-objects object-printer inferior-producer 
                                    \key duplicate-key duplicate-test}

This function is responsible for generating all of the graph node output records
of the graph.  \arg{graph-record} is the graph output record, and \arg{stream}
is the output stream.  The graph node output records are generating by calling
the object printer on the root objects, then (recursively) calling the inferior
producer on the root objects and calling the object printer on all inferiors.
After all of the graph node output records have been generated, the value of
\cl{graph-root-nodes} of \arg{graph-record} must be set to be a sequence of the
those graph node output records that correspond to the root objects.

\arg{root-objects}, \arg{object-printer}, \arg{inferior-producer},
\arg{duplicate-key}, and \arg{duplicate-test} are as for
\cl{format-graph-from-roots}.

\Defgeneric {layout-graph-nodes} {graph-record stream arc-drawer arc-drawing-options} 

This function is responsible for laying out the nodes in the graph contained in
the output record \arg{graph-record}.  It is called after the graph output has
been collected, but before the graph record has been displayed.  The method on
\cl{standard-graph-output-record} implements the usual graph layout constraint
solver.  \arg{stream} is the stream on which the graph is displayed.

\Defgeneric {layout-graph-edges} {graph-record stream arc-drawer arc-drawing-options}

This function is responsible for laying out the edges in the graph.  It is
called after the graph nodes have been layed out, but before the graph record
has been displayed.  The method on \cl{standard-graph-output-record} simply
causes thin lines to be drawn from each node to all of its children.
\arg{graph-record} and \arg{stream} are as for \cl{layout-graph-nodes}.


\Defprotoclass {graph-node-output-record}

The protocol class that represents a node in graph; a subclass of
\cl{output-record}.
\IfYouWantClass {a} {graph node output record} {graph-node-output-record}

\Defpredicate {graph-node-output-record-p} {object}

Returns \term{true} if \arg{object} is a \term{graph node output record},
otherwise returns \term{false}.

\Defclass {standard-graph-node-output-record}

The instantiable class of output record that represents a graph node.  Its
parent will be a graph output record.  This is a subclass of
\cl{graph-node-output-record}.

\Defgeneric {graph-node-parents} {graph-node-record}

Returns a sequence of the graph node output records whose objects are
``parents'' of the object corresponding to the graph node output record
\arg{graph-node-record}.  Note that this is not the same as
\cl{output-record-parent}, since \cl{graph-node-parents} can return output
records that are not the parent records of \arg{graph-node-record}.

\Defgeneric {(setf graph-node-parents)} {parents graph-node-record}

Sets the parents of \arg{graph-node-record} to be \arg{parents}.  \arg{parents}
must be a list of graph node records.

\Defgeneric {graph-node-children} {graph-node-record}

Returns a sequence of the graph node output records whose objects are
``children'' of the object corresponding to the graph node output record
\arg{graph-node-record}.  Note that this is not the same as
\cl{output-record-children}, since \cl{graph-node-children} can return output
records that are not child records of \arg{graph-node-record}.

\Defgeneric {(setf graph-node-children)} {children graph-node-record}

Sets the children of \arg{graph-node-record} to be \arg{children}.
\arg{children} must be a list of graph node records.

\Defgeneric {graph-node-object} {graph-node-record}

Returns the object that corresponds to the graph node output record
\arg{graph-node-record}.  It is permissible for this function to work correctly
only while inside of the call to \cl{format-graph-from-roots}.  It is unspecified
what result will be returned outside of \cl{format-graph-from-roots}.  This
restriction is permitted so that CLIM is not required to capture application
objects that might have dynamic extent.
